\begin{abstract}
Indoor environmental management requires more than displaying temperature and relative humidity (RH); residents must decide when to dehumidify, ventilate, go out, or dry laundry, and these decisions depend on outdoor conditions at a different place and time. Indoor sensor readings and outdoor weather forecasts differ in spatial coverage, temporal validity, and data semantics, and different appliance operations can produce similar indoor trajectories. This paper presents an evidence-traceable, domain-specific agent that integrates single-room sensing, township-level forecasts, versioned domain rules, and user feedback through constrained tool routing and rule-guided reasoning. An ESP32 node publishes measurements over MQTT to an SQLite store; a Python analytics engine performs quality checks, trend analysis, and event detection; and a large language model (LLM) is restricted to intent parsing, tool selection, and explanation, with every recommendation bound to an evidence record of values, sources, validity windows, and rule versions. An event-triggered feedback loop with event aggregation and a query cooldown collects time-aligned operation labels while bounding user interruptions, and keeps unanswered queries as unknown rather than negative. We define a two-week single-room evaluation protocol with event-grouped temporal splits that compares a fixed-rule baseline with a decision tree, two feedback policies, and template versus agent responses in terms of macro-F1, interaction burden, evidence consistency, and latency. Measured deployment results are not yet available; the protocol and target metrics defined here form the planned evaluation benchmark.
\end{abstract}

\begin{IEEEkeywords}
context-aware decision support, domain-specific agents, environmental sensing, human-in-the-loop feedback, Internet of Things
\end{IEEEkeywords}
